\section{Nuevo paradigma: Environment + Task-Reward}

Esta sección cierra con la perspectiva planteada en el programa: el núcleo de los datos podría pasar de las anotaciones input-output a la construcción de environment y task-reward. Es decir, en el futuro el entrenamiento de un Agent no consistirá solo en recopilar respuestas, sino en construir entornos interactivos, definir tareas y diseñar reward verificables para que el modelo haga self-improve a partir de la experiencia.

\begin{figure}[H]
\centering
\includegraphics[width=0.96\textwidth]{new-agent-paradigm.png}
\caption{Nuevo paradigma de Agent: el núcleo de los datos pasa de input-output a environment + task-reward. Diagrama conceptual propio, elaborado a partir de la conversación de 02:06:06--02:15:20.}
\end{figure}

\begin{knowledgebox}{Lectura de la figura: de datos de respuestas a datos de experiencia}
Los datos input-output enseñan al modelo a «ver una entrada y dar una respuesta»; environment + task-reward le enseñan a «actuar en un entorno y mejorar a partir de la retroalimentación». Esta es la diferencia de paradigma de datos entre un Agent y un modelo conversacional común.
\end{knowledgebox}

\subsection{Rubrics as reward y self-improvement}

La sección anterior trató el nuevo paradigma de datos; esta se centra en su problema central: de dónde sale la reward. El programa menciona rubrics as reward, es decir, usar criterios de evaluación como mecanismo de recompensa. Su importancia radica en que muchas tareas reales no tienen un simple correcto o incorrecto, pero sí pueden tener criterios de evaluación por niveles. Que un Agent pueda mejorarse a sí mismo depende de que logre encontrar o construir una reward verificable y aprovechar de forma eficaz la experiencia de interacción.

\begin{warningbox}{Condiciones previas de la self-improve}
La automejora de un Agent no significa dejar que el modelo se autocomplazca sin límite. Requiere tareas verificables, entornos confiables, trayectorias que admitan auditoría y mecanismos de seguridad que prevengan el reward hacking.
\end{warningbox}

\begin{knowledgebox}{Tabla comparativa de los nuevos paradigmas}
\begin{tabularx}{\textwidth}{p{0.22\textwidth}p{0.34\textwidth}X}
\toprule
Paradigma & Forma de los datos & Cuello de botella principal \\
\midrule
Input-output & Entradas y respuestas de referencia & Costo de anotación y límites de generalización. \\
Trajectory & Observaciones, acciones, llamadas a herramientas, resultados & Calidad de las trayectorias y recuperación de errores. \\
Environment & Mundo de tareas interactivo & Costo de construcción del entorno y cobertura. \\
Task-reward & Definición de tareas y retroalimentación verificable & Diseño de la reward y seguridad. \\
Experience reuse & Automejora a partir de la experiencia histórica & Memoria, eliminación de ruido y prevención de la autocomplacencia. \\
\bottomrule
\end{tabularx}
\end{knowledgebox}

\subsection{Family of Agents}

La reward y la self-improvement tratan de cómo aprende un sistema individual; esta sección pasa a la organización de múltiples Agent. Al final se menciona que el Agent es como un «cerebro extendido» respaldado por un ejército. La metáfora indica que el Agent del futuro quizá no sea una sola entidad, sino un conjunto de agentes especializados: coding, search, browser, memory, safety, entre otros, que colaboran entre sí.

\begin{figure}[H]
\centering
\includegraphics[width=0.92\textwidth]{family-of-agents.png}
\caption{Family of Agents: el Agent es como un cerebro extendido, respaldado por un ejército. Diagrama conceptual propio, elaborado a partir de la conversación de 02:15:20--02:16:41.}
\end{figure}

\subsection{Resumen de la sección}

El nuevo paradigma de Agent integra datos, entornos, recompensas y ciclos de experiencia. La verdadera ingeniería de sistemas consiste en lograr que estos ciclos sean escalables, verificables y desplegables de forma segura.
